% Chapter introduction: why the topic arose, who first asked, what it answers.
% STATUS: written 2026-11-05 (cites only the chapter's own sources).
The theme of this chapter is a single observation that grew into a method.
Posner showed in 1957 that a prime ring admitting a nonzero centralizing
derivation must be commutative
\cite[p.~385]{Bresar1993}. Mayne proved the parallel statement
for centralizing automorphisms \cite[p.~385]{Bresar1993}.
Bre\v{s}ar then asked what an arbitrary additive centralizing map looks
like \cite[p.~385]{Bresar1993}. He answered it in the prime
case: such a map is multiplication by a central element, up to a central
remainder \cite[Theorem~A, p.~385]{Bresar1993}. Section
\ref{sec:functional-identities:commuting-maps} reports this line and its
one-sided-ideal refinement of Lee
\cite[Theorem~2, p.~335]{LeeTK1997Centralizing}.

Out of these commuting-map theorems grew the general theory.
Section~\ref{sec:functional-identities:fi-theory} introduces a functional
identity as an equation with functions as unknowns, and the $d$-free set
as the concept on which the whole development rests
\cite[p.~9 and p.~x]{BresarMartindaleChebotar2007FI}.
Section~\ref{sec:functional-identities:d-free} reports the threefold
development of that concept from the preface of the monograph
\cite[p.~x]{BresarMartindaleChebotar2007FI}.
Section~\ref{sec:functional-identities:applications-to-derivations}
carries the theory back to derivations and shows how a two-variable
functional identity settles Jordan $\sigma$-derivations on prime rings
\cite[Theorem~2.1, p.~221]{LeeTK2015JordanFI}.
Section~\ref{sec:functional-identities:centralizers-multipliers} closes
with the twisted form: a map that commutes with an endomorphism still has
the central form on semiprime rings
\cite[Theorem~2, p.~2948]{LeeTK2001SigmaCommuting}.

\hbfigure{09-functional-identities/commuting}{Commuting maps on prime rings (Bre\v{s}ar): the only additive maps commuting with every element are the obvious ones.}{fig:fi:commuting}
